(1) 位置算符 \(x=\sqrt{\hbar/(2m\omega)}(a+a^\dagger)\) 仅连接量子数相差1的 Fock 态。初态 n=0，非零的激发矩阵元只有
\[
x_{10}=\langle1|x|0\rangle=\sqrt{\hbar/(2m\omega)}.
\]
对 n≥2，\(\langle n|x|0\rangle=0\)。

(2) 交互绘景的一阶跃迁振幅
\[
c_1^{(1)}(t)=-\frac i\hbar\int_0^t V_{10}e^{i\omega t'}dt'
=\frac{iFx_{10}}{\hbar}\frac{e^{i\omega t}-1}{i\omega}
=\frac{Fx_{10}}{\hbar\omega}(e^{i\omega t}-1).
\]
这是对任意有限 t 的一阶小 F 结果。

(3) 取模平方得到
\[
P_{0\to1}(t)\simeq |c_1^{(1)}|^2
=\frac{4F^2x_{10}^2}{\hbar^2\omega^2}\sin^2\frac{\omega t}2
=\frac{2F^2}{m\hbar\omega^3}\sin^2\frac{\omega t}2.
\]
概率是二阶于 F，因为一阶振幅取模平方；若此估计不再远小于1，则微扰展开失效。

(4) \(t\to0\) 时 \(\sin(\omega t/2)\sim\omega t/2\)，故 \(P\simeq F^2t^2/(2m\hbar\omega)\)。当 \(t=2\pi/\omega\) 时一阶跃迁振幅相消，概率到该阶为零；真实系统更高阶项仍可能贡献，不能把微扰零点称为严格无跃迁。
